\clearpage
\thispagestyle{fancy}
\begin{center}
    {\Large \textbf{Arabic Summary}}
\end{center}


\begin{arabtext}
\setcode{utf8}
تُوسِّع هذه الرسالة خليّة تجميع روبوتية تعاونية ثنائية الأذرع ـــ مبنية حول مُحرِّك صناعي (\L{ABB IRB 120}) وذراع بحثية (\L{AR4 Mk3}) تتشاركان طاولة عمل واحدة ـــ من نموذج تسلسلي للالتقاط والوضع إلى منظومة \textbf{متوازية التنفيذ}، \textbf{واعية بالسلامة}، \textbf{وقابلة للتشغيل عن بُعد}. ويأتي هذا العمل بوصفه المرحلة الثانية من المشروع، حيث تُشكِّل خطوط الإدراك والمعايرة واستنتاج القبضة والتوجيه البصري المُنجزة في المرحلة الأولى الأساس الذي بُنيت عليه \textbf{خمس مساهمات} جديدة.

يستبدل هذا العمل آلة الحالة أحادية الخيط السابقة بمُشرف غير متزامن متعدد الخيوط مبني على نظام (\L{ROS 2 Jazzy Jalisco}) يُرسل أهداف الحركة إلى الذراعين في آنٍ واحد، مع عزل مُراقب السلامة في مجموعة استدعاءات مستقلة. ويجمع خط تخطيط هجين بين (\L{MoveIt Task Constructor}) واستدعاءات (\L{MoveIt 2}) لتخطيط حركات الذراع الواحدة وعملية التسليم المُنسَّقة ذات الاثنتي عشرة درجة حرية. ويُصنِّف مدير اكتشاف المناطق المخصَّص كل ذراع ضمن مناطق عمل حصرية أو مشتركة، ويُجهض الأهداف غير الآمنة عبر مسار مقاطعة يقل عن المللي ثانية. كما يوفِّر التوأم الرقمي متعدد الأنماط القائم على (\L{Gazebo}) مرآةً حيّةً للحركة الفيزيائية، وبيئة اختبار آمنة تُنفَّذ فيها البرامج المحفوظة قبل أن يتلقّى أيُّ متحكم حقيقي هدفاً. وتُكوِّن بنية الاتصال الصناعي ثنائية القناة ـــ (\L{Fast DDS}) عبر شبكة (\L{ZeroTier}) للموضوعات الفورية و(\L{Syncthing}) لملفات البرامج ـــ جسراً بين حاسوبين عبر الإنترنت العامة، مع قناة برمجية مخصَّصة للإيقاف الطارئ عن بُعد.

أسفر التحقق التجريبي عن انخفاض زمن دورة الالتقاط والوضع بنسبة \L{\textbf{46.5\%}} وزمن خمول كل ذراع بنسبة \L{\textbf{86.0\%}} مقارنةً بمرجعية المرحلة الأولى. وحافظ مدير اكتشاف المناطق على فصل أدنى قدره \L{14.5 cm} بين المؤثرات النهائية للذراعين، بمتوسط زمن استجابة \L{0.089 ms}؛ كما تتبَّع التوأم الرقمي الحركة الحيّة بفارق بصري دون \L{50 ms}، ورفض أي برنامج غير مُتحقَّق منه؛ وانتقل أمر الإيقاف الطارئ من طرف إلى طرف في أقل من \L{5 ms}. وتُختتم الرسالة بتحليل الجدوى الصناعية الذي يُغطي الجدوى الاقتصادية، والتكامل التقني، والتوافق مع معياري السلامة (\L{ISO 10218}) و(\L{ISO/TS 15066})، والاعتبارات البشرية للمُشغِّل عن بُعد.
\end{arabtext}
